% from ega4-20-fr.tex
\oldpage[IV-4]{250}
de $X$ dans $Y$ relativement à $S$ sont équivalents au sens de \sref{IV.20.2.1}, d'où une application canonique
\[
\label{IV.20.5.2.1}
  \operatorname{Ps.hom}_{X/S}(X,Y)\to\operatorname{Ps.hom}_S(X,Y).
  \tag{20.5.2.1}
\]
\end{env}

\begin{proposition}[20.5.3]
\label{IV.20.5.3}
Supposons $Y$ séparé sur $S$.
Alors :
\begin{enumerate}[label=\emph{(\roman*)}]
  \item L'application \sref{IV.20.5.2.1} est injective et identifie $\operatorname{Ps.hom}_{X/S}(X,Y)$ au sous-ensemble de $\operatorname{Ps.hom}_S(X,Y)$ formé des pseudo-$S$-morphismes $\omega$ tels que $\operatorname{dom}_S(\omega)$ soit universellement schématiquement dense relativement à $S$.
  \item Le préfaisceau $U\mapsto\operatorname{Ps.hom}_{U/S}(U,Y)$ sur $X$ est un faisceau.
\end{enumerate}
\end{proposition}

\begin{proof}
L'assertion~(i) est immédiate, car dire que deux $S$-morphismes $u:U\to Y$, $u':U'\to Y$ sont équivalents au sens de \sref{IV.20.2.1} signifie alors que ce sont les restrictions d'un même morphisme $\operatorname{dom}_S(\omega)\to Y$ \sref{IV.20.2.4}, et si $U$ et $U'$ sont universellement schématiquement denses par rapport à $S$, il en est de même \emph{a fortiori} de $\operatorname{dom}_S(\omega)$.
Pour prouver~(ii), notons que $U\mapsto\operatorname{Ps.hom}_S(U,Y)$ est alors un faisceau \sref{IV.20.2.6} ; d'autre part, étant donné un recouvrement ouvert $(U_\alpha)$ de $U$, et un pseudo-$S$-morphisme $\omega$ de $U$ dans $Y$, pour que $\operatorname{dom}_S(\omega)$ soit universellement schématiquement dense dans $U$ relativement à $S$, il résulte aussitôt des définitions (cf.\ \sref{IV.20.2.1}) qu'il faut et il suffit que chacun des ensembles $\operatorname{dom}_S(\omega)\cap U_\alpha=\operatorname{dom}_S(\omega|U_\alpha)$ soit universellement schématiquement dense dans $U_\alpha$ relativement à $S$ ; d'où~(ii).
\end{proof}

Lorsque $Y$ est séparé, on notera $\mathscr{P}\!s.\!\operatorname{hom}_{X/S}(X,Y)$ le sous-faisceau
\[
  U\mapsto\operatorname{Ps.hom}_{U/S}(U,Y)
\]
de $\mathscr{P}\!s.\!\operatorname{hom}_S(X,Y)$.

Lorsque $X$ est \emph{plat et localement de présentation finie} sur $S$ et $Y$ séparé sur $S$, les pseudo-morphismes de $X$ dans $Y$ relativement à $S$ s'identifient encore aux pseudo-$S$-morphismes $\omega$ tels que, pour toute fibre $X_s$ du morphisme $X\to S$, $\operatorname{dom}_S(\omega)\cap X_s$ soit \emph{schématiquement dense} dans $X_s$ \sref{IV.11.10.10}.

\begin{env}[20.5.4]
\EGAReaderTarget{ega4.c04.statement.20-5-4}{IV.20.5.4}
Un cas particulièrement important où $Y$ est séparé sur $S$ est le cas où $Y=S[T]=S\otimes_{\bb{Z}}\Spec(\bb{Z}[T])$ ($T$ indéterminée).
Il y a alors correspondance biunivoque entre les pseudo-$S$-morphismes de $X$ dans $Y$ et les \emph{pseudo-fonctions} sur $X$ \sref{IV.20.2.8} en vertu de la définition d'un produit de $\bb{Z}$-préschémas.
Les pseudo-morphismes de $X$ dans $Y$ relativement à $S$ s'identifient alors, en vertu de \sref{IV.20.5.3}, aux \emph{pseudo-fonctions $\varphi$ sur $X$ telles que $\operatorname{dom}(\varphi)$ soit universellement schématiquement dense relativement à $S$}.
Le faisceau $\mathscr{P}\!s.\!\operatorname{hom}_{X/S}(X,Y)$ est un \emph{sous-faisceau d'anneaux} de $\sh{M}'_X$, que l'on note $\sh{M}'_{X/S}$.

On pose alors $\operatorname{Ps.hom}_{X/S}(X,Y)=M'(X/S)$ et on dit que ses éléments sont les \emph{pseudo-fonctions sur $X$ relatives à $S$}.
\end{env}

\begin{env}[20.5.5]
\label{IV.20.5.5}
Soient $X$, $Y$, $Z$ trois $S$-préschémas, $f:Y\to Z$ un $S$-morphisme.
On peut, dans le raisonnement de \sref{IV.20.3.1}, remplacer partout « schématiquement dense » par « universellement schématiquement dense relativement à $S$ » ; pour tout pseudo-morphisme $\omega\in\operatorname{Ps.hom}_{X/S}(X,Y)$, les morphismes $f\circ u$, où $u$ parcourt $\omega$, appartiennent donc à une même classe d'équivalence (pour la relation définie dans \sref{IV.20.5.1}), que l'on
